\chapter{Profiling, performance and hardware backends}
\label{evod47:chapter}

An optimization changes a kernel from 100 to 80 ms. Its speedup is 1.25, while
latency falls 20\%. Now suppose that kernel occupied only one fifth of the
request's critical path. Even an ideal twofold kernel improvement would reduce
the request to 90\% of its old duration. Optimizing the wrong region can make
an impressive profiler screenshot irrelevant to the user.

This chapter retains the package's top-down measurement ladder and synthetic
request trace. It adds explicit time populations, failure denominators,
critical-path arithmetic, uncertainty and backend validation contracts.
The executable analysis runs on assigned numbers. It reports no measured
CPU, Metal, CUDA or Vulkan throughput and makes no hardware ranking.

\section{Correctness is a precondition for comparison}
\index{correctness parity}
\label{evod47:correct}
Fix model/artifact, tokenizer, position convention, accepted input stream,
state lifetime, output semantics and numerical tolerance before timing.
For Chapter 31's restricted native recurrence, compare hidden state and logits
with its scalar path; do not call it the full DOGMA architecture. For Chapter
35 use the actual Chapter 23 learned-absolute-position multi-layer Hermon
Transformer and its original full causal oracle, including every layer's K/V.
A detached rotary one-head fixture would test a different model.

Test multiple shapes, empty/partial inputs, masks, capacity boundaries,
permutations and cancellation. Close logits do not guarantee the same selected
token near a tie. If a lower-precision or approximate kernel changes sampling
behavior, report numerical and task-quality effects under the new contract.
Correctness tests are necessary; they do not establish learned capability.

\section{Freeze a workload and a measurement population}
\index{benchmark record}\index{open-loop load}\index{closed-loop load}\index{warmup}
\label{evod47:workload}
A benchmark record includes commit/artifact and data hashes, hardware/firmware,
OS, driver, compiler, backend libraries, precision/quantization, scheduler,
memory budget, power/clocks, input/output lengths, arrival process, concurrency,
cache policy, sampling, warmup, repetitions and raw outcomes. Preserve offered,
accepted, rejected, cancelled, failed and completed counts separately.
Report peak memory and energy only when actually measured. For energy per
accepted token, declare the power-sampling boundary, idle subtraction, integration
interval and accepted-token denominator; device and whole-host energy differ.

An open-loop generator schedules arrivals independently of response time;
a closed-loop generator waits for completion or a concurrency slot. Both are
useful, but a slow closed-loop service reduces offered load automatically.
It can hide the overload a production arrival stream would create. Fix and
replay traces, include burstiness, and state whether the load generator itself
became the bottleneck. Batching a preconstructed full tensor does not evaluate
continuous serving under arrivals.

Separate cold artifact loading, compilation/graph capture, allocator startup
and cache warming from steady state unless cold start is the target metric.
Specify warmup by iterations or a declared stationarity rule, and archive
discarded warmup samples. Do not choose the cutoff after inspecting the best
performance region. Thermal/power changes can invalidate stationarity.

\section{Trace timestamps into user-visible metrics}
\index{TTFT}\index{ITL}\index{TPOT}\index{Goodput}\index{survivor bias}
\label{evod47:metrics}
For request $r$, use a single observation clock: arrival $a_r$, first output
token $t_{r,1}$ and subsequent token times $t_{r,2},\ldots,t_{r,n_r}$ in ms.
The first token is not repeated in the subsequent-times tuple. Define
\[
 TTFT_r=t_{r,1}-a_r,\qquad
 ITL_{r,i}=t_{r,i}-t_{r,i-1}\quad(i\ge2).
\]
There is no ITL for a one-token response. A failed zero-output request has no
TTFT sample, not a zero-millisecond success. It requires an explicit terminal
timestamp; requests with outputs may also specify one after their last token.
Keep failures in the outcome denominator and check terminal time against the window.
Cross-host clocks need synchronization and error bounds before subtraction.

\Plate{01-timestamps}{The three supplied synthetic requests have different arrivals and token gaps. The first token counts once; TTFT includes waiting before it. The long third request violates both the 100 ms TTFT and 25 ms maximum-gap targets.}{evod47:timestamps}
\Listing{Request.validate}

For a closed completed cohort observed over $[s,e]$, raw output-token rate is
$\sum_r n_r/((e-s)/1000)$ in tokens/s. Our goodput counts successful requests
meeting TTFT and \emph{every} ITL bound, divided by the same seconds.
This maximum-gap SLO is stricter than an average time-per-output-token
(TPOT) target. A request with gaps $(1,99)$ ms has mean 50 ms but fails a
60 ms maximum-gap objective. Do not compare systems using different SLOs
under the same goodput label. DistServe uses a phase-sensitive TTFT/TPOT
framework~\cite{distserve47}; we do not reproduce its experiment.

\Listing{metrics}
The supplied trace has arrivals $(0,10,20)$, first tokens $(40,70,150)$ and
subsequent tokens $(55,70,85)$, $(90,110,130)$, $(180,210)$ ms. Over 250 ms,
there are 11 tokens: 44 tokens/s. Two of three requests meet the 100/25 ms
TTFT/ITL bounds: 8 good requests/s and attainment $2/3$. TTFTs are $(40,60,130)$;
nearest-rank p50/p95 are 60/130 ms and pooled p95 ITL is 30 ms. Every number
is synthetic arithmetic, not a service measurement.

The analyzer requires all arrivals and emitted tokens to lie in its declared
window. It is a closed-cohort analysis, not a streaming estimator with pending
requests. A real offered-load SLO must also count requests unresolved at the
deadline. Silently dropping them creates survivor bias. Token-weighted pooled
ITL and request-weighted worst-gap distributions are different populations;
long outputs contribute more samples to the first.

\section{Percentiles and uncertainty need declared estimators}
\index{percentile}\index{uncertainty}\index{paired runs}
\label{evod47:uncertainty}
For $n$ sorted durations and $0<q\le1$, nearest rank selects
$x_{(\lceil qn\rceil)}$. The reference rejects empty samples, nonfinite or
negative times and invalid quantiles. Missing ITL is represented as no value,
not zero. A four-sample p95 is the maximum and cannot characterize a rare
production tail reliably.
\Listing{percentile}

Repeat independent runs with matched workload traces and randomized execution
order where feasible. Report sample sizes, run-level dispersion, raw samples
and exclusions. Thousands of correlated token gaps from one run are not
thousands of independent experiments. Block bootstrap or a hierarchical
model can respect trace/run clustering, but assumptions and interval method
must be explicit. The fixture computes only mean, sample standard deviation
and range of matched-run speedup ratios; it supplies no confidence interval.
\Listing{paired_ratio}
For paired durations $(100,120,80)$ and $(80,100,64)$ ms, ratios are
$(1.25,1.2,1.25)$. Their mean is about 1.233 and sample deviation 0.0289.
This is different from dividing arbitrary unpaired best runs. Repetitions
that share a thermal drift or queue regime are not independent evidence.

\section{Measure a critical path before summing durations}
\index{Critical path}\index{Amdahl model}
\label{evod47:path}
A latency decomposition such as queue + prefill + transfer + decode + tools
+ serialization is exact only for nonoverlapping serial phases covering the
whole request. If retrieval overlaps prefill, summing both overcounts.
Represent a dependency DAG with span duration $d_i$ and predecessor set
$pred(i)$; assuming sufficient resources and earliest-start scheduling,
\[
 E_i=d_i+\max_{j\in pred(i)}E_j,\qquad
 L=\max_i E_i,
\]
where the maximum of an empty predecessor set is zero. Resource contention
must be represented by additional edges or measured waiting spans.
\Listing{critical_path}

The exact fixture has queue 10 ms, then prefill 40 and retrieval 30 ms in
parallel, followed by decode 20 ms waiting for both. The path is 70 ms,
whereas summing all spans gives 100. Halving retrieval to 15 changes no
critical-path latency. Halving prefill to 20 gives $10+30+20=60$ ms; retrieval
becomes the limiting branch.

\Plate{02-timeline}{Parallel spans are not additive latency components. After shortening prefill, the retrieval branch becomes critical. Solid connectors denote dependencies; horizontal extent denotes assigned duration.}{evod47:timeline}

For a serial fraction $f$ sped up by $s$, Amdahl's model gives
\[
 S_{total}=\frac{1}{(1-f)+f/s}.
\]
Derive it by dividing original total time $T$ into $(1-f)T$ unchanged and
$fT/s$ improved. With $f=0.2,s=2$, total speedup is $1/0.9\approx1.111$,
even though the region doubled. It assumes fixed work and no new overhead,
queue feedback or overlap changes; service optimization often violates those
assumptions. Measure the resulting system, not only the local ratio.

\section{Count arithmetic and traffic at a named boundary}
\index{Arithmetic intensity}\index{roofline}\index{memory traffic}
\label{evod47:roofline}
For $F$ operations and $B$ bytes moved at a specified memory level,
arithmetic intensity $I=F/B$ has units FLOP/byte. With compute ceiling $P$
FLOP/s and bandwidth $W$ bytes/s, a throughput ceiling is
\[
 P_{attainable}\le\min(P,WI),\qquad I^*=P/W.
\]
Equivalently, ideal execution time is at least $\max(F/P,B/W)$ under these
resource assumptions. Adding synchronization/launch overhead can only increase
time. The Roofline model provides useful ceilings~\cite{roofline47}; peaks,
measured sustainable bandwidth and actual traffic must not be interchanged.
\Listing{roofline_bound}

\Plate{03-roofline}{Assigned ceilings of 100 TFLOP/s and 2 TB/s decimal meet at 50 FLOP/byte. The point at intensity 2 has a 4 TFLOP/s ceiling. These curves are a model, not a device profile.}{evod47:rooflineplate}

The same operation can have different intensity at register, shared/on-chip,
last-level-cache and device-memory boundaries. Tensor element counts alone
do not reveal all device traffic. Include output writes, staging, spills,
alignment, metadata and parameter reads where the scope requires them.
Comparing useful operations to total physical bytes can differ sharply from
comparing nominal dense operations to ideal tensor payloads.

\subsection{Attention and native examples without a false speed ranking}
For one decode query per layer/head, QK and AV interactions cost approximately
$4LHtd_h$ multiply/add FLOPs, excluding softmax, projections and other layers.
Reading K/V once costs $2LHtd_hb$ logical bytes. Their ratio is $2/b$ FLOP/byte.
At two-byte elements, it is one, independent of retained length; total work
and bytes still grow with $t$. Batch reuse, GQA, cache layout and kernel fusion
change the traffic. FlashAttention's IO-aware tiling is a relevant mechanism
comparison~\cite{flash47}; it does not make this CPU oracle an optimized kernel.

Chapter 34's actual restricted recurrence counts $2BDV+4BD$ multiply/add
operations, sigmoid evaluations separately, and resident payload
$B(3D+V)b$. Gather/scatter adds $4BDb$. With batch $B=8$, width $D=64$,
vocabulary $V=4$ and element width $b=8$, the generated values are:
\begin{center}\input{\Generated/costs.tex}\end{center}
The attention row uses Chapter 35's large storage geometry, not the small
two-layer executed fixture. The native operation count excludes sigmoid and
shared parameter traffic; it is not a complete FLOP count. These rows cannot
rank capabilities, latency or energy. They expose declared accounting scopes.

\section{Profile top down with tools appropriate to the backend}
\index{profiling}\index{CPU backend}\index{Metal backend}\index{CUDA backend}\index{Vulkan backend}
\label{evod47:backends}
First measure offered-load SLOs and request phases. Then inspect system
timeline gaps, transfers, allocator stalls and synchronization. Only after
locating a critical region inspect operators and kernel counters. PyTorch
Profiler can collect operator activity, shapes/stacks and device activity
where supported~\cite{torch47}. Detailed collection itself adds overhead;
confirm final performance with an appropriately uninstrumented protocol.

Nsight Systems supports CPU/GPU system tracing and focused regions; Nsight
Compute analyzes CUDA kernels~\cite{nsys47,ncu47}. Use phase annotations to
connect prefill, decode, native steps, memory and transfer to a request trace.
Counter collection/replay can perturb cache state and timing. Architecture-specific
occupancy is evidence to interpret, not a universal objective to maximize.

\begin{center}\small
\begin{tabular}{p{19mm}p{54mm}p{65mm}}
\toprule Backend & Record explicitly & Investigate on that platform\\\midrule
CPU & ISA, core/NUMA topology, affinity, threads & Vectorization, caches, memory bandwidth, dispatch\\
Metal & SoC, OS/framework, precision, memory pressure & Native GPU capture/counters, scheduling, bandwidth\\
CUDA & GPU, driver/toolkit, clocks/power, sharing & System timeline, kernel metrics, interconnect\\
Vulkan & Device/driver, shader compiler, subgroup features & Barriers, timestamp support, layout, vendor counters\\
\bottomrule
\end{tabular}
\end{center}

These are experiment fields, not measurements. Apple's Metal tooling supports
platform GPU performance investigation~\cite{metal47}; unified memory does
not remove bandwidth/capacity limits or synchronization. Vulkan exposes
explicit synchronization responsibilities~\cite{vulkan47}; portability of an
API does not imply portable performance. Do not mechanically translate CUDA
counters into Metal or Vulkan terms. Unsupported operators, precision and
fallbacks must be disclosed before comparing backends.

\section{Time asynchronous work at an intentional boundary}
\index{synchronization}\index{device timing}
\label{evod47:sync}
Host timing around an enqueue can measure submission rather than completion.
Device events or a synchronized wall-clock boundary are needed for completed
work, with exact scope documented. Synchronizing every operation can destroy
the overlap the service normally uses. Compare equivalent timing boundaries
and include required output materialization. CPU/GPU clock domains and profiler
timestamps are not automatically interchangeable.

Batching may improve device occupancy but increase queue delay and token gaps.
Sweep concurrency and arrival rate, report occupancy distributions, and
distinguish useful accepted tokens from padded/cancelled work. For distributed
placement, measure queue, serialization, network, reconstruction and commit
with overlap and clock error stated. Chapter 46's bandwidth model is a
hypothesis to calibrate against real transfer traces, not a measured input.

\section{An optimization loop that can reject its own hypothesis}
\label{evod47:loop}
Freeze baseline and workload; reproduce correctness; locate the critical
region; state one mechanism and predicted signature; change it; recheck
correctness; remeasure paired runs. For example, fusion predicts fewer
intermediate bytes and launches. It may fail because register pressure or
specialization overhead grows. Archive that failure, raw traces and all
configuration changes. A speedup without the predicted signature invites a
different causal explanation such as changed workload or warmup.

The reference tests exact synthetic outcomes, percentile permutation
invariance, invalid clocks, missing-token outcomes, maximum-gap versus mean
SLOs, a throughput/goodput reversal, DAG overlap, roofline units, Amdahl and
the earlier chapters' declared traffic arithmetic. It launches no hardware
backend. Backend quality/performance experiments against completed Chapters
34/38/46 remain integration work; synthetic success is never their substitute.

\section{Terms and the application boundary}
\Term{Goodput} counts successful work satisfying a declared service objective.
\Term{Critical path} is the longest causal duration under stated resource
assumptions. \Term{Arithmetic intensity} names both work and a memory boundary.
The industrial contracts now define what a service would have to preserve
and measure. They do not establish production readiness. Chapter 48 returns
to genomic applications, where task-specific biological evidence remains a
separate requirement from a correct or fast runtime.

\section{Twelve worked exercises}
\begin{enumerate}
\setlength{\itemsep}{1pt}\setlength{\parsep}{0pt}\setlength{\parskip}{0pt}
\item \textbf{Fastest kernel, fastest service?} \emph{Solution:} only if its
critical contribution and system effects support that conclusion. The parallel
retrieval example changes no latency when shortened alone.
\item \textbf{Compute TTFT and ITL.} Arrival 10, first 70, then 90 and 110 ms.
\emph{Solution:} TTFT 60 ms; gaps 20,20 ms. First token is counted once.
\item \textbf{Throughput versus goodput.} Why retain failed requests?
\emph{Solution:} omitting them inflates attainment. Goodput counts successful
SLO-qualified requests per common observation time; raw tokens can include
partial failed responses if explicitly reported.
\item \textbf{Calculate a ceiling.} $I=2$ FLOP/B and $W=2$ TB/s.
\emph{Solution:} bandwidth ceiling 4 TFLOP/s decimal; take the minimum with
compute peak. It is not achieved throughput.
\item \textbf{Name the byte boundary.} Is tensor size equal to DRAM traffic?
\emph{Solution:} no; caches, rereads, writes, spills, staging and metadata alter it.
\item \textbf{Warmup honestly.} Remove the slowest first ten samples after seeing
results? \emph{Solution:} not as a predeclared steady-state protocol. State and
archive the rule before collecting comparison evidence.
\item \textbf{Trace before counters.} What might dominate besides kernels?
\emph{Solution:} queueing, copies, launch gaps, allocator waits and synchronization.
Use the critical path to choose detailed profiling scope.
\item \textbf{Quote CUDA peak.} Does it establish application performance?
\emph{Solution:} no; collect workload timings and relevant achieved counters,
with exact precision, sharing mode, versions and clocks.
\item \textbf{Unified memory.} Does Metal avoid all movement cost?
\emph{Solution:} no; bandwidth, capacity, locality and synchronization still matter.
Use native evidence rather than a CUDA analogy.
\item \textbf{Speedup arithmetic.} 20 to 10 ms is $2\times$; 100 to 80 is
$1.25\times$ and 20\% reduction. \emph{Solution:} divide old time by new for
speedup, divide saved time by old time for fractional reduction.
\item \textbf{Synthetic versus measured.} May the 44 tokens/s trace be labeled
a GPU benchmark? \emph{Solution:} no; all timestamps were assigned. A hardware
claim needs archived real samples and the complete measurement contract.
\item \textbf{Design a backend comparison.} \emph{Rubric:} fix model/tolerance,
workload/arrivals and versions/hardware; separate warmup; time completed work;
archive repeated raw outcomes, tails/goodput/memory, explanatory profiles and
uncertainty. Disclose unsupported paths and failed changes.
\end{enumerate}

